Between 1950 and 1975 millions of Spaniards left the villages of the interior for Madrid, Barcelona and other cities, and in Spain people still talk about going back to the village on holiday. I wanted to see whether this shows in where people actually sleep. The Spanish transport ministry publishes mobile-phone estimates of overnight stays, and from them I build one directed network per night: the nodes are 2,435 zones of mainland Spain (municipalities or groups of small ones), and a link goes from zone $i$ to zone $j$ when at least 10 residents of $i$ sleep in $j$. I compare an ordinary weekday in February 2024 with the summer holiday of 14--15 August and with two family holidays, All Saints and Christmas Eve, and I set these networks against census data on where each city's residents were born. On the summer night the undirected network is almost four times denser and splits into sending zones, whose residents sleep in many other zones, and receiving zones, which host people from many. But random networks with the same degrees reproduce most of the change in how its nodes mix, so most of it comes from the degrees alone. On every night, a city sends more people to the other provinces where a larger share of its residents were born. Whether this pull grows in August depends on the night I compare with: it grows against October weekdays, but I cannot detect it against February. It does grow clearly for residents who moved to their current province before 1981, the exodus generation, and for those born in villages, and, against February, only for stays inland. In the fifth of zones that lost the largest share of their 1900--1960 peak population, 24\,\% receive more people than they send on the ordinary night and 99\,\% on 15 August. The data link cities to provinces, not travellers to their own villages, so I cannot say that people go back to their village; but the inland part of the summer fits the map of the exodus.
